\noindent\textbf{摘 要}

本文针对华北山区乡村在 2024--2030 年的农作物种植策略问题，围绕有限耕地、多季种植、作物适宜性、轮作连作约束以及未来市场不确定性，建立了由确定性优化、随机规划和相关性扩展组成的种植决策模型。模型以地块、作物、季次和年份为核心索引，将附件中的地块类型、作物适宜规则、2023 年种植记录和亩产成本价格数据转化为统一的种植可行集合，并以 2023 年各作物分季总产量作为基准预期销售量。

对于问题一，本文建立多周期混合整数线性规划模型，以七年总利润最大为目标，同时刻画地块面积上限、最小种植比例、同地块连作禁止、三年内豆类轮作和作物分散度等约束。针对超产滞销和超产半价销售两种情形分别求解，得到滞销情形下七年利润为 \textbf{3552.16 万元}的时限内可行方案，半价销售情形下七年利润为 \textbf{6110.96 万元}。其中，滞销情形受混合整数求解时间限制，论文中仅作为可行解与拓扑固定线性规划精化结果表述；半价情形在多组配置下利润波动仅为 0.031\%，表现出较好的求解稳定性。

对于问题二，本文在问题一半价方案的种植拓扑基础上，引入需求、成本、价格、亩产波动和露天灾害等随机因素，构建基于样本平均逼近的两阶段随机规划模型，并用 95\% CVaR 描述尾部风险。模型在 $K=50$ 个情景下转化为拓扑固定线性规划求解，得到七年期望利润 \textbf{6429.27 万元}，较问题一半价方案提高约 318.31 万元；情景总利润标准差为 112.03 万元，最差 5\% 情景下平均仍保持盈利，说明该方案在收益提升的同时具备一定抗风险能力。

对于问题三，本文进一步分析作物之间的统计关系，采用 Spearman 相关系数刻画替代关系，并通过 K-means 聚类识别作物功能分组；在此基础上引入同簇负相关作物的需求转移机制和前茬豆类带来的互补性成本折减机制。扩展模型求得七年期望利润 \textbf{6443.76 万元}，较问题二提高 14.49 万元，同时总成本由 906.02 万元降至 891.60 万元，表明替代与互补关系主要通过优化作物结构和降低后茬成本改善方案表现。

为检验模型可靠性，本文从维度一致性、手算特例、随机情景稳定性、求解收敛性和模板回填一致性等方面进行验证，35 项一致性测试全部通过；灵敏度分析显示，风险厌恶系数 $\lambda$ 和最小种植比例 $\delta$ 对利润影响较小，灾害概率是主要风险来源，互补折减系数 $\phi$ 在 0.85--0.95 区间内不改变主要结论。综合来看，本文模型能够在满足农作制度和销售约束的前提下，给出兼顾收益、风险和作物结构协调性的七年种植方案。

\vspace{0.8em}
\noindent\textbf{关键词：} 农作物种植策略 \quad 混合整数线性规划 \quad 两阶段随机规划 \quad CVaR \quad 样本平均逼近 \quad 作物相关性
\newpage
